\documentclass[12pt]{article}
\usepackage[utf8]{inputenc}
\usepackage[T1]{fontenc}
\usepackage{hyperref}
\title{Totonac, Yecuatla Language Resources for K-12 Curriculum Development}
\author{Compiled by Claude Sonnet 5 (claude-sonnet-5) for Andr\'as Kornai}
\date{}
\begin{document}
\maketitle

\section{Introduction}
Yecuatla Totonac (ISO 639-3 tlc) is an endangered Totonacan language of Veracruz, Mexico. The Endangered Languages Project lists it under the name Misantla Totonac, with Yecuatla, San Marcos Atexquilapan and Jilotepec as dialects. Very little material was found.

\section{Primary Educational Resources}
The Endangered Languages Project resources page mentions a 2010 Totonac learning resource, \emph{Aprendamos totonaco}. It was not verified whether this covers the Yecuatla variety.

\section{Linguistic and Reference Materials}
\begin{itemize}
\item Glottolog entry and bibliography: \url{https://glottolog.org/resource/languoid/id/yecu1235}
\item Endangered Languages Project resources: \url{https://endangeredlanguages.com/node/38959/resources}
\end{itemize}

\section{Religious and Community Resources}
None found.

\section{Digital and Audio Resources}
Global Recordings Network has a page for this language, but it reports that no recordings are currently available: \url{https://globalrecordings.net/en/language/tlc}

\section{Curriculum Development Context}
Few materials exist for this variety, and it is spoken mainly by older adults. Any curriculum would depend on documentation and recording work with remaining speakers.
\end{document}